\section{第 3 章 系统需求分析}

\subsection{3.1 总体需求概述}

本系统面向的并不是单一类型用户，而是同时覆盖普通天气查询用户、具有一定专业分析需求的用户以及农业应用场景下的业务型用户。因此，系统需求既包括基础的信息展示需求，也包括更深层的分析需求和一定程度的场景化决策支持需求。为了避免系统在功能组织上出现分散和割裂，本文将“预报分析能力”作为主线，在此基础上扩展历史统计、探空诊断、误差校正和农业服务等模块。

\subsection{3.2 用户需求分析}

普通用户最根本的需求就是得到迅速、清楚、可信、连续的天气信息。因此系统需要有易于理解的页面入口、当前天气概览、未来数小时变化趋势、未来数天变化趋势，在视觉上尽可能做到信息集中、切换方便、内容直观。

对于具备一定专业需求的用户而言，单一模式的天气结果往往不足以支持判断。他们更希望看到不同模式之间的对比关系、预报起报时次的合理解释、上空层结结构与稳定度参数，以及历史上同类型天气在多年样本中的相对异常程度。这类需求决定了系统不能只停留在“数据展示”，而要具备一定分析和解释能力。

对于农业场景用户而言，仅仅知道未来降水或温度并不够，还需要将天气信息与作物生长过程、积温进度和农业风险联系起来。因此，系统还应具备面向作物的积温估算、灾害风险提示和农事建议能力。

\subsection{3.3 功能需求分析}

根据目前系统的实现情况，本文把功能需求归纳为以下几方面。

第一，系统要具备多模式预报分析功能。该功能包含72小时精细化逐小时预报、多模式对比分析、预报趋势展示和ECMWF预报机器学习校正结果接入等，是整个平台的主要功能。

第二，系统要有探空分析功能。此功能可以对指定的时次进行探空数据抓取，得到层结表解析结果，绘制 Skew-T 和参数图，计算 CAPE、CIN、K 等稳定度参数。当抓取失败、时次不一致或者参数异常的时候，也应当具有一定的回退和异常提示功能。

第三，系统应该有历史气候分析的功能。该功能需要具备年度页、趋势页的展示能力，可以给出气温、降水等指标的长期统计结果，根据原始极值和EW指数分别代表事件两个方面来分析某一年里具有代表性的异常天气过程。

第四，系统要具有农业气象服务功能。该功能要依靠作物知识库，给典型作物提供积温进度、成熟日期估算、农业气象风险提醒以及农事建议。

第五，系统要有导航以及综合展示的功能。首页导航页要担负起平台总入口、主要模块分发、实时天气概览以及智能提醒汇总等任务，让用户在同一个页面上先建立起对于当前天气和系统能力结构的总体认识，然后进入相应的专业模块进行查看。

\subsection{3.4 非功能需求分析}

在非功能需求方面，系统至少需要满足以下要求。

响应性为第一。因为系统包含了很多外部数据源，如果没有缓存、并行请求或者失败回退等机制，那么页面加载就会出现明显迟滞。因此系统要依靠缓存、并发抓取以及接口分层来保证主要页面的可用性。

第二，稳定性。系统外部数据源有开放预报接口、探空资料源、历史观测数据源，不同数据源响应格式、可用时次、异常情况各不相同，所以系统要具有较强异常处理、容错能力。

第三，具有可扩展性。系统后续会增加更多的站点、更多的预报模式、更多的作物或者更复杂的机器学习模型，因此在架构设计的时候不能过于耦合，尽量使用分模块的方式。

第四，可以维护。由于系统包含页面模板、接口逻辑、数据抓取、分析计算、模型加载等多种代码，所以要保证模块的职责尽量清晰，便于以后的迭代。

\subsection{3.5 业务主线界定}

为了防止系统变成单纯的项目功能清单，本文在需求分析阶段就将预报分析能力的建立确定为业务主线。多模式预报模块直接给出未来天气分析能力，探空分析模块加强了专业的诊断能力，历史气候模块给目前的年份以及某些天气过程赋予了长时间的背景支撑，机器学习误差校正模块对预报品质展开后处理改良，农业气象模块在此基础上扩展到面向行业应用。这样一种组织方式既可以保证系统的综合性，又不会造成各个模块之间的相互隔离。

